\documentclass[11pt]{article}
\usepackage{fontspec}
\usepackage[french]{babel}
\usepackage{amsmath}
\usepackage{amssymb}
\usepackage{amsthm}
\usepackage{mathrsfs}
\usepackage{tikz-cd}
\usepackage{enumitem}

\newcommand{\printedpage}[1]{\marginpar{\footnotesize p.~#1}}
\newcommand{\msp}[1]{\marginpar{\footnotesize [#1]}}
\newcommand{\Ch}{\underline{\mathrm{Ch}}}
\newcommand{\CH}{\underline{\mathrm{CH}}}
\newcommand{\Ty}{\underline{\mathrm{Ty}}}
\newcommand{\Hom}{\underline{\mathrm{Hom}}}
\newcommand{\red}{\mathrm{r\acute{e}d}}

\begin{document}
\setcounter{footnote}{20}

\printedpage{49}
% head: A. Grothendieck : Dérivateurs, Ch. VII, Édition des Universités de Montpellier II et Paris VII
\msp{page 74}
des flèches inversibles de $X'$ (lequel ensemble satisfait bien aux deux conditions d'admissibilité). Donc par ce qui précède on a un diagramme commutatif
\[
\begin{tikzcd}
X \ar[r,"\kappa"] \ar[d,"f"'] & \Pi_{\Sigma}(X) \ar[d,"\bar f = \Pi_{\Sigma}(f)"] \ar[dl,dotted,"\tilde f"'] \\
X' \ar[r,"\kappa'","\sim"'] & \Pi_{\Sigma'}(X'),
\end{tikzcd}
\]
où $\bar f$ est uniquement déterminé. Comme $\kappa'$ est un isomorphisme (comme on a vu tantôt, puisque $\Sigma'$ \texttt{[est]} formé des flèches inversibles), l'existence et l'unicité de $\bar f$ équivaut à la conclusion annoncée,

\hfill c.q.f.d.

\textbf{Encore un sorite} sur les classes d'homotopie d'éléments d'une catégorie $X$. De façon générale, on voit que $X$ et $\Pi_{\Sigma}(X) = X\Sigma^{-1}$ ont mêmes composantes connexes, car si
\msp{page 75}
$a$ et $b$ dans $X$ sont reliés par une flèche de $X\Sigma^{-1}$, i.e. par un $\Sigma$-chemin, alors ils sont dans une même composante connexe de $X$. Mais l'inverse est également vrai, comme on voit p.~ex. en prenant $\Sigma = \mathrm{Fl}(X)$, et en disant que dans le \textit{groupoïde} $\Pi(X)$, où deux sommets sont dans une même composante \texttt{[connexe]} si et seulement si $\exists\ x \longrightarrow y$ (relation qui est non seulement transitive et reflexive, mais aussi symétrique dans un groupoïde). Un chemin
\[
a \xrightarrow{\quad c \quad} b
\]
sera appelé aussi une \textit{homotopie} de $a$ vers $b$. On dit que $a$ et $b$ sont \textit{homotopes} s'il existe une homotopie $c$ de $a$ vers $b$, auquel cas il existe aussi une homotopie `inverse' $c^{o}$ de $b$ vers $a$. Cela signifie aussi que $a$ et $b$ sont dans la même composante connexe de $X$, i.e. définissent le même élément de $\pi_0(X)$.

L'homotopie sera dite \textit{simple} si $c$ est de type $\Delta_{+}$ ou $\Delta_{-}$, i.e. si elle est donnée par une flèche $u : a \longrightarrow b$ ou $v : b \longrightarrow a$. Si on cherche une homotopie entre $a$ et $b$, on essaiera de la trouver la plus `courte' possible. Si on a une homotopie $c$ de longueur $l$ de $a$ vers $b$, alors en remplaçant $\tau$ par $\tau_{\red}$, donc $I_{\tau}$ par $I_{\tau\,\red}$, le sous-intervalle réduit
\msp{page 76}
correspondant, on trouve une homotopie de $a$ à $b$ de type $\tau_{\red}$, dont la longueur est égale à la \textit{`longueur utile'} du type $\tau$ de $c$ (ce qui justifie la terminologie). Notons que si
\[
\tau = \Delta_{\varepsilon}^{l_1}\Delta_{-\varepsilon}^{l_2}\Delta_{\varepsilon}^{l_3}\cdots\Delta_{(-1)^{r}\varepsilon}^{l_r},
\]
alors
\[
\tau_{\red} = \Delta_{\varepsilon}\Delta_{-\varepsilon}\Delta_{\varepsilon}\cdots\Delta_{(-1)^{r}\varepsilon},
\]
% foot: Version 22 décembre 2001
\printedpage{50}
% head: A. Grothendieck : Dérivateurs, Ch. VII, Édition des Universités de Montpellier II et Paris VII
\texttt{[plutôt deux fois} $(-1)^{r-1}$\texttt{]} où $r$ est la longueur utile de $\tau$ (la longueur pas utile étant $n = \sum\limits_{i\in[1,r]} l_i$).

Quand on part d'une homotopie entre \textit{foncteurs}
\[
X \longrightarrow Y,
\]
il est entendu qu'il s'agit d'homotopies dans la catégorie $\Hom(X,Y)$. (Ou, suivant le contexte, dans une sous-catégorie précisée de $\Hom(X,Y)$.)

\msp{page 77}
À la page 61, j'ai invoqué l'existence d'un \textit{foncteur} de \textit{`composition des chemins'}
\[
\Ch_{\Sigma}(X;a,b)\times\Ch_{\Sigma}(X;b,c) \ \longrightarrow\ \Ch_{\Sigma}(X;a,c),
\]
qui sur la fibre du premier membre au dessus de $(\tau,\tau') \in \operatorname{Ob}(\Ty\times\Ty)$, serait donné par
\[
\Ch_{\Sigma,\tau}(X;a,b)\times\Ch_{\Sigma,\tau'}(X;b,c) \ \longrightarrow\ \Ch_{\Sigma,\tau\tau'}(X;a,c), \tag*{$(*)$}
\]
la composition des chemins de types précisés $\tau$, $\tau'$. Qu'on ait là un \textit{foncteur,} pas seulement une application sur les ensembles d'objets, est assez évident, et peut s'exprimer par exemple à l'aide du diagramme d'amalga\texttt{[ma]}tion
\[
\begin{tikzcd}[column sep=small]
 & \varepsilon = \{0,1\} \ar[dl,hook'] \ar[dr,hook] & \\
I_{\tau} \ar[dr,hook] & & I_{\tau'} \ar[dl,hook'] \\
 & I_{\tau\tau'} &
\end{tikzcd}
\]
en le transformant par $\Hom(\cdot,X)$ (au lieu de le transformer par $\mathrm{Hom}_{\mathrm{Cat}}(\cdot,X)$ seulement), du moins dans le cas particulier où $\Sigma = \mathrm{Fl}(X)$ (qui était celui de la page 61). Mais le passage aux flèches est assez évident - pour deux composés $c'_1\circ c_1$ et $c'_2\circ c_2$, de
\[
\begin{tikzcd}[column sep=small]
a \ar[rr,"c_1"] \ar[d,equal] & {} \ar[d,"w"'] & b \ar[rr,"c'_1"] \ar[d,equal] & {} \ar[d,"w'"] & c \ar[d,equal] \\
a \ar[rr,"c_2"] & {} & b \ar[rr,"c'_2"] & {} & c
\end{tikzcd}
\]
\texttt{[on]} voit que les homomorphismes $w''$ du premier dans le second, dans $\Ch_{\Sigma}(X;a,c)$, qui sont l'identité sur le sommet d'indice $n+1$ des deux chemins (où $n = \operatorname{long}\tau$), correspondant à $1_{\tau}$ et à $0_{\tau'}$, `sont' exactement les couples
% foot: Version 22 décembre 2001
\printedpage{51}
% head: A. Grothendieck : Dérivateurs, Ch. VII, Édition des Universités de Montpellier II et Paris VII
\msp{page 78}
formés d'un homomorphisme $w : c_1 \longrightarrow c_2$ et $w' : c'_1 \longrightarrow c'_2$ dans $\Ch_{\Sigma,\tau}(a,b)$ et $\Ch_{\Sigma,\tau'}(b,c)$ respectivement, et si $w''$, $w$, $w'$ se correspondent ainsi, on posera
\[
w'' = w * w'.
\]
Il y a une transitivité évidente à vérifier pour voir que ça définit un \textit{foncteur} $(*)$. On peut exprimer la situation à l'aide du diagramme
\[
\begin{tikzcd}[column sep=small]
 & \Ty\times\Ty \ar[rrr,"\pi\ (\text{composition})"] \ar[dl,"p_1"'] \ar[dr,"p_2"] & & & \Ty \\
\Ty & & \Ty, & &
\end{tikzcd}
\]
où $p_1$, $p_2$ sont les deux projections. On considère sur $\Ty\times\Ty$ les trois catégories fibrés % [sic] fibrés -> fibrées, agreement with "catégories"
\[
p_1^*(\CH_{\Sigma,\tau}(a,b)),\ \ p_2^*(\CH_{\Sigma,\tau'}(b,c)),\ \ \pi^*(\Ch_{\Sigma,\tau\tau'}(a,c))
\]
\texttt{[plutôt} $\pi^*(\CH_{\Sigma,\tau\tau'}(a,c)$\texttt{]}, dont les fibres en $(\tau,\tau')$ sont respectivement % [sic] \pi^*(\CH_{\Sigma,\tau\tau'}(a,c) -> \pi^*(\CH_{\Sigma,\tau\tau'}(a,c)), one closing parenthesis missing
\[
\Ch_{\Sigma,\tau}(a,b),\ \ \Ch_{\Sigma,\tau'}(b,c),\ \ \Ch_{\Sigma,\tau\tau'}(a,c).
\]
(On écrit $\CH$ au lieu de $\Ch$, pour indiquer que l'on considère les catégories envisagées comme des catégories (fibrées) \textit{au dessus de} $(\Ty)$ (\footnote{Mais il s'agit pour $p_1^*$, $p_2^*$ et $\pi^*$ de produits cartésiens ordinaires dans Cat, il n'y a pas à avoir peur et à écrire des $\CH$ \ldots}).) Au dessus de $\Ty\times\Ty$, on a un homomorphisme (cartésien) de catégories fibrées
\[
p_1^*(\CH_{\tau}(a,b))\times_{\Ty\times\Ty} p_2^*(\CH_{\tau}(b,c)) \ \longrightarrow\ \pi^*(\CH_{\tau}(a,c)) \tag*{$(**)$}
\]
\texttt{[plutôt} $p_1^*(\CH_{\tau}(a,b))\times_{\Ty\times\Ty} p_2^*(\CH_{\tau'}(b,c)) \ \longrightarrow\ \pi^*(\CH_{\tau\tau'}(a,c))$\texttt{]},
\msp{page 79}
car c'est OK sur les fibres ($(*)$ page 77), et si
\[
(\varphi,\varphi') : (\tau_1,\tau'_1) \longrightarrow (\tau,\tau') \qquad (\varphi : \tau_1 \longrightarrow \tau,\ \varphi' : \tau'_1 \longrightarrow \tau')
\]
est une flèche dans $\Ty\times\Ty$, on en déduit un diagramme commutatif
\[
\begin{tikzcd}
\Ch_{\Sigma,\tau}(a,b)\times\Ch_{\Sigma,\tau'}(b,c) \ar[r,"\text{comp.}"] \ar[d,"\varphi^*_{a,b}\times(\varphi'_{b,c})^*"] & \Ch_{\Sigma,\tau}(a,c) \ar[d,"(\varphi*\varphi')^*_{a,c}"] \\
\Ch_{\Sigma,\tau_1}(a,b)\times\Ch_{\Sigma,\tau'_1}(b,c) \ar[r,"\text{comp.}"] & \Ch_{\Sigma,\tau_1\tau'_1}(a,c)
\end{tikzcd}
\]
% foot: Version 22 décembre 2001
\printedpage{52}
% head: A. Grothendieck : Dérivateurs, Ch. VII, Édition des Universités de Montpellier II et Paris VII
\texttt{[plutôt} $\Ch_{\Sigma,\tau\tau'}(a,c)$\texttt{]}, i.e. la composition des chemins d'extrémités fixés est compatible avec les images inverses de chemins par flèches entre les types de chemins.

Mais l'homomorphisme cartésien $(**)$ de catégories fibrées sur $\Ty\times\Ty$ n'est pas tout à fait la même chose qu'un anodin foncteur
\[
\Ch_{\Sigma}(a,b)\times\Ch_{\Sigma}(b,c) \ \longrightarrow\ \Ch_{\Sigma}(a,c). \tag*{$(*)$}
\]
La formule $(*)$ définit en tout cas une application sur les ensembles d'objets, puisque les ensembles d'objets de chacune des trois catégories envisagées se décompose en somme d'ensembles $\Ch_{\Sigma,\tau}(\bullet,\bullet)$, où $\tau$ parcourt $\mathrm{Ty} = \operatorname{Ob}\Ty$. Mais cela se prolonge-t-il de façon naturelle aux morphismes?

\msp{page 80}
Mais voilà que je me rappelle que ce fameux produit cartésien $p_1^*(\text{fourbi})\times p_2^*(\text{machin})$ n'est autre que le produit cartésien tout bête
\[
\Ch_{\Sigma}(a,b)\times\Ch_{\Sigma}(b,c),
\]
envoyé sur $\Ty\times\Ty$ par le produit cartésien des deux foncteurs de projection
\[
\begin{aligned}
p_{a,b} : \Ch_{\Sigma}(a,b) &\ \longrightarrow\ \ \Ty\\
p_{b,c} : \Ch_{\Sigma}(b,c) &\ \longrightarrow\ \ \Ty,
\end{aligned}
\]
d'où
\[
\begin{tikzcd}
\Ch_{\Sigma}(a,b)\times\Ch_{\Sigma}(b,c) \ar[d,"p_{a,b}\times p_{b,c}"] \\
\Ty\times\Ty
\end{tikzcd}
\]
\msp{page 81}
qui est bel et bien fibrant (comme produit cartésien de deux foncteurs fibrants) et a les fibres qu'il faut! On trouve donc bien, par le sorite général sur les foncteurs fibrants (ici sur $\Ty\times\Ty$) un foncteur de catégories fibrées sur $\Ty\times\Ty$
\[
\Ch_{\Sigma}(a,b)\times\Ch_{\Sigma}(b,c) \ \longrightarrow\ \qquad \underbrace{\pi^*(\CH_{\Sigma}(a,c))}_{=\,(\Ty\times\Ty)\times_{\Ty}\Ch_{\Sigma}(a,c)\ \ \text{(\footnotemark)}} \qquad .
\]
\footnotetext{produit fibré par $\pi : \Ty\times\Ty \longrightarrow \Ty$.}
Comme on a d'autre part un foncteur du dernier membre (produit fibré) vers le facteur $\Ch_{\Sigma}(a,c)$, on gagne. (On l'aurait retrouvé par le raisonnement direct que je m'apprêtais de faire, et que j'ai biffé - mais je suis plus heureux comme ça!)

% foot: Version 22 décembre 2001
\printedpage{53}
% head: A. Grothendieck : Dérivateurs, Ch. VII, Édition des Universités de Montpellier II et Paris VII
Pour bien faire, il faudrait prouver l'\textit{associativité} des foncteurs de composition, pas seulement au niveau des objets, par des arguments comme ceci, mais ça revient simplement à une relation du type
\[
(w''*w')*w = w''*(w'*w)
\]
pour un système de trois flèches dans
\[
\Ch_{\Sigma}(a,b),\ \Ch_{\Sigma}(b,c)\ \ \text{et}\ \ \Ch_{\Sigma}(c,d)
\]
respectivement. Ça ne peut être que trivialement
\msp{page 82}
vrai, mais fastidieux à expliciter.

Les sueurs froides sont terminées!

En résumé, le formalisme des $\Ch_{\Sigma}(X)$ et des $\Ch_{\Sigma}(X;a,b)$ développé jusqu'à présent donne satisfaction complète sur les points suivants :

\begin{enumerate}
\item[a)] Bon formalisme de composition des chemins (qui était dégueulasse avant).

\item[b)] Traitement élégant de la construction des catégories de fractions $\mathcal{M}\Sigma^{-1}$. Mais je m'attends que ce sera moins joli qu'avant, quand il s'agira de remplacer un $(\mathcal{M},\Sigma)$ par des $(\ \underbrace{\mathcal{M}(X)}_{=\,\Hom(X,\mathcal{M})}\ ,\Sigma_X)$ ($X$ une catégorie d'indices) \texttt{[plutôt} $\Hom(X^{o},\mathcal{M})$\texttt{]}. C'est pour y pallier qu'il me faudra développer des énoncés de comparaison, entre la catégorie $\Ch_{\Sigma}$ construite ici, et d'autres dont le type dans Hot se prête aux calculs de façon plus directe, comme dans l'état premier de la question. De plus, on sera sûrement obligé d'en venir à des Ind-catégories, pour pouvoir suivre la variance de $\mathcal{M}(X)\Sigma_X^{-1}$ pour $X$ variable.

\item[c)] Bon formalisme des `homotopies' et des compositions des homotopies (ce qui est une autre façon de parler des chemins et de leur composition).
\end{enumerate}
% foot: Version 22 décembre 2001

\end{document}
